% =========================
% Chapter 4
% =========================
\chapter{System Design}
\label{ch:system_design}

\section{From Proposal to Final Design}
\label{sec:design_evolution}

The original project proposal envisioned a purpose-built wearable device --- camera,
microphone, and speaker in a dedicated enclosure --- streaming to cloud services. Two
realizations during development reshaped this design into its final, stronger form.

First, \textbf{the smartphone already is the wearable}. A modern phone integrates
every sensor the custom device would have carried --- camera, microphone, speaker,
accelerometer, magnetometer --- in hardware that is mass-produced, familiar to the
user, and already in their pocket. Building custom hardware would have reproduced
this integration work at prototype quality while adding battery, enclosure, and
connectivity problems that contribute nothing to the research questions.

Second, \textbf{the browser removes the last installation barrier}. Delivering the
client as a plain web page --- no app store, no APK, no platform approval --- means a
user (or an evaluator) points their phone at a URL and the system works. The
requirements this imposes (HTTPS for camera access, iOS audio unlock quirks,
permission gestures) proved solvable, and are documented in
Section~\ref{sec:frontend}.

The navigation task evolved even more substantially. The proposal described
\emph{collaborative waypoint identification}: the user names a sequence of landmarks,
and the system walks them landmark to landmark. Early experimentation exposed the
flaw: in an unfamiliar building, a blind user cannot name intermediate landmarks ---
that is precisely the information they lack, and demanding it contradicts the
cognitive-load motivation of the whole project. The final design inverts the
responsibility: the user names only the \emph{destination}, and the system explores
--- scanning each room, recognizing arrival semantically, choosing doors, and
guiding the user through them. The user contributes what they genuinely have (coarse
direction knowledge, confirmation of transit) rather than what they lack. What
remains from the proposal is the deeper principle both designs share: perception
scoped to an explicit task, under explicit state control.

\section{Requirements}
\label{sec:requirements}

\subsection{Functional Requirements}
\begin{itemize}[nosep]
  \item \textbf{FR1} --- Voice-initiated tasks: the user starts, completes, and
        aborts every task by voice (push-to-talk), with UI buttons as fallback.
  \item \textbf{FR2} --- Object Allocation: locate a requested object from a curated
        COCO vocabulary and guide the user to it with direction and distance cues.
  \item \textbf{FR3} --- Reach Guidance: when the target is within arm's reach and a
        hand is visible, guide the hand to the target and complete the task
        automatically on touch.
  \item \textbf{FR4} --- Navigation: guide the user to a named destination room
        without maps or infrastructure, recognizing arrival from room-indicator
        objects and traversing doors as needed.
  \item \textbf{FR5} --- Obstacle warnings: while the user walks during navigation,
        warn about objects and unnamed clutter in the walking lane, preempting all
        other guidance.
  \item \textbf{FR6} --- Information queries: ``describe'' the current view and
        ``repeat'' the last utterance on demand during a task.
  \item \textbf{FR7} --- Session continuity: an interrupted connection (dropped
        Wi-Fi, page reload) resumes the active task on reconnect.
\end{itemize}

\subsection{Non-Functional Requirements}
\begin{itemize}[nosep]
  \item \textbf{NFR1} --- Zero install: the client is a web page; the only
        client-side requirement is a modern browser.
  \item \textbf{NFR2} --- Responsiveness: guidance must follow scene changes closely
        enough to feel live (detection at ${\sim}3$--$5$\,FPS; spoken reaction
        within ${\sim}1$--$2$\,s).
  \item \textbf{NFR3} --- Speech discipline: no utterance may be cut off by another;
        ambient guidance must never queue up and lag behind reality.
  \item \textbf{NFR4} --- Trustworthy perception: nothing is spoken to the user
        unless it survived temporal-consistency and (for doors) semantic
        verification; false guidance is treated as a safety defect, not a cosmetic
        one.
  \item \textbf{NFR5} --- Testability: all decision logic must be executable without
        cameras, models, or network, so behavior is verifiable in an automated
        suite.
  \item \textbf{NFR6} --- Cane-complementary: the system provides advance
        information; it does not replace the white cane's ground-truth contact zone.
\end{itemize}

\section{Architecture Overview}
\label{sec:architecture}

Lumen is a thin-client / smart-server system connected by a single WebSocket
(Figure~\ref{fig:architecture}). The browser client captures and forwards sensor
data and plays back audio; it holds no task state. The backend owns everything else:
speech recognition, command parsing, the session state machine, the per-task
perception loops, and speech synthesis. Serving the frontend from the same process
(same origin) means one tunnel or host exposes the entire system.

\begin{figure}[H]
\centering
\resizebox{0.98\textwidth}{!}{%
\begin{tikzpicture}[node distance=6mm]
  % Client side
  \node[box, minimum width=44mm] (cam) {getUserMedia\\camera 640$\times$480};
  \node[box, minimum width=44mm, below=of cam] (ptt) {MediaRecorder\\push-to-talk audio};
  \node[box, minimum width=44mm, below=of ptt] (sens) {Motion + compass\\(IMU, magnetometer)};
  \node[box, minimum width=44mm, below=of sens] (audio) {Audio queue\\primed \texttt{<audio>} (iOS)};
  \node[draw, rounded corners=3pt, fit=(cam)(ptt)(sens)(audio), inner sep=4mm,
        label={[font=\small\bfseries]above:Browser client (phone)}] (client) {};

  % Server side
  \node[box, minimum width=50mm, right=34mm of cam] (router) {Router + Session\\(one per connection)};
  \node[box, minimum width=50mm, below=of router] (fsm) {Task FSM\\(authoritative state)};
  \node[box, minimum width=50mm, below=of fsm] (tasks) {Task engines\\Object Allocation / Reach\\Navigation (exploration)};
  \node[box, minimum width=50mm, below=of tasks] (svc) {Services: faster-whisper,\\parser, YOLOv8n/m, door models,\\MediaPipe, SegFormer, gTTS};
  \node[draw, rounded corners=3pt, fit=(router)(fsm)(tasks)(svc), inner sep=4mm,
        label={[font=\small\bfseries]above:FastAPI backend (single process)}] (server) {};

  % WS link
  \draw[arr, <->] (client.east) -- node[lbl, above]{WebSocket}
         (server.west);
\end{tikzpicture}}
\caption[System architecture]{System architecture: a sensing-only browser client and a single-process
backend joined by one WebSocket carrying JSON control messages and tagged binary
media.}
\label{fig:architecture}
\end{figure}

\subsection{Wire Protocol}
\label{sec:protocol}

One WebSocket carries two channels. \emph{Text frames} hold JSON control messages:
lifecycle events (\texttt{user\_event}: start / stop / confirm), device sensing
(\texttt{motion}: still / moving / walking; \texttt{heading}: compass degrees,
throttled to ${\geq}2^\circ$ change at ${\leq}5$\,Hz), session resumption
(\texttt{resume}), and, server-to-client, state pushes (\texttt{fsm\_state}),
transcriptions, haptic cues, and errors. \emph{Binary frames} carry media, each
prefixed by a one-byte tag: \texttt{0x01} JPEG camera frames (client$\to$server,
${\sim}5$\,FPS), \texttt{0x02} WebM/Opus push-to-talk recordings
(client$\to$server), and \texttt{0x03} MP3 speech clips (server$\to$client). The
protocol is frozen and documented, and unknown messages are logged and ignored,
which lets client and server evolve independently.

\subsection{Session Model}

Each WebSocket connection owns a \texttt{Session}: the FSM instance, the latest
decoded frame and its timestamp, the latest motion state and compass heading, the
task context (target object or destination), and a handle to the running task loop.
There is no shared task state across users --- multi-user isolation is a property of
the architecture. A resume pool preserves the active task for five minutes across
disconnects: when the phone drops Wi-Fi mid-task and reconnects, the client presents
its previous session id and the server re-fires the task with a spoken
``Resuming\ldots'' announcement (FR7).

\section{The Task Finite State Machine}
\label{sec:fsm}

\begin{figure}[H]
\centering
\resizebox{0.98\textwidth}{!}{%
\begin{tikzpicture}[node distance=16mm and 34mm]
  \node[state] (idle) {Idle};
  \node[state, right=of idle] (listen) {Listening\\ForCommand};
  \node[state, above right=of listen] (obj) {ObjectAllocation\\Active};
  \node[state, below right=of listen] (nav) {Navigation\\Active};
  \node[state, right=72mm of listen] (ret) {Returning\\ToIdle};

  \draw[arr] (idle) -- node[lbl, above]{user\_start} (listen);
  \draw[arr] (listen) -- node[lbl, above left]{command\_recognized\\(find X)} (obj);
  \draw[arr] (listen) -- node[lbl, below left]{command\_recognized\\(navigate to Y)} (nav);
  \draw[arr] (obj) -- node[lbl, above right]{task\_complete /\\user\_stop} (ret);
  \draw[arr] (nav) -- node[lbl, below right]{task\_complete /\\user\_stop} (ret);
  \draw[arr] (obj.west) to[bend right=12] node[lbl, above]{task\_abort} (listen.north);
  \draw[arr] (nav.west) to[bend left=12] node[lbl, below]{task\_abort} (listen.south);
  \draw[arr] (ret.south) to[bend left=22] node[lbl, below]{cleanup\_done} (idle.south);
\end{tikzpicture}}
\caption[Session finite state machine]{The session FSM. Every transition is driven by an explicit event; the only
autonomous transitions are \texttt{cleanup\_done} (one tick after entering
ReturningToIdle) and \texttt{task\_complete} fired by a task engine on detected
success (grasp, arrival).}
\label{fig:fsm}
\end{figure}

The FSM is authoritative: the client renders whatever state the server pushes and
maintains no state of its own. Entry and exit hooks on the two active states start
and cancel the corresponding task engine's asynchronous loop, so a task can never
outlive its state --- cancellation, completion, disconnection, and errors all funnel
through the same teardown path. Reach Guidance runs as a phase \emph{inside}
ObjectAllocationActive rather than as a separate FSM state, because the transition
between far-field guidance and hand guidance is a perception decision (target near
and hand visible), not a user decision.

\section{Voice Interaction Design}
\label{sec:voice_design}

Push-to-talk, rather than an always-open microphone, keeps the user in control and
sidesteps both privacy concerns and false activations. Utterances flow through
faster-whisper, then the fuzzy command parser
(Section~\ref{sec:command_parsing}), which classifies them as task commands
(``find my cup'', ``navigate to the kitchen''), completions (``got it'' /
``stop''), or information queries (``describe'', ``repeat that''). The parser's
navigation vocabulary is deliberately kept in lockstep with the set of rooms the
navigation engine can actually recognize --- a destination the parser accepts but
the explorer cannot find would produce two contradictory spoken responses.

Every accepted command is confirmed by voice (``Looking for your cup.''), every
rejection prompts a retry (``I didn't catch that, please repeat.''), and every task
end is announced. Haptic cues complement voice at key moments --- entering reach
range, touching the target, task completion --- degrading silently on devices
without a vibration API.

\section{Safety Design}
\label{sec:safety_design}

Four principles govern safety-relevant behavior:

\begin{enumerate}[nosep]
  \item \textbf{Never speak unverified perception.} Objects require temporal
        consistency (3-of-5 frames in Object Allocation; localized sighting counts
        in Navigation); doors additionally pass geometric gates and semantic
        verification (Section~\ref{sec:door_funnel}).
  \item \textbf{Obstacles preempt everything.} While the user is walking, the
        obstacle watchdog's warnings override door guidance until the path clears.
  \item \textbf{The camera is the odometer.} ``You walked through the door'' is
        only ever inferred when sustained camera motion accompanies the visual
        evidence --- detection flicker alone can never fake progress.
  \item \textbf{Bounded autonomy.} Exploration checks in with the user after a
        configurable number of rooms (``say stop anytime\ldots''), and a
        whole-journey timeout guarantees the system never wanders indefinitely;
        ``stop'' is honored from any state at any time.
\end{enumerate}

Lumen is designed as a complement to the white cane: the cane owns the contact zone
(${\sim}1$\,m); Lumen's value is the 2--4\,m advance information --- where the door
is, what is in the path, which room this is. This division of labor is communicated
to users and maintained throughout the guidance phrasing, which directs attention
and movement but never promises collision safety.
